\documentclass[12pt]{article}
\usepackage[utf8]{inputenc}
\usepackage[T1]{fontenc}
\usepackage{hyperref}
\title{Yerukula Language Resources for K-12 Curriculum Development}
\author{Compiled by Claude Sonnet 5.5 (claude-sonnet-5-5) for Andr\'as Kornai}
\date{}
\begin{document}
\maketitle

\section{Introduction}
Yerukula (also Yerukala, Kurru basha, Kulavatha; ISO 639-3 \texttt{yeu}, Glottocode yeru1240) is a South Dravidian language of India. Speaker figures differ: the 2011 census gives 58,065 native speakers, while Omniglot says about 70,000. This survey found no school materials and no dedicated dictionary. Little is available online, and this report lists only what the searches showed.

\section{Primary Educational Resources}
No primer, textbook or curriculum for Yerukula was found in search results. Omniglot says that a script devised by Professor Prasanna Sree exists, and that another script was used in a number of school textbooks. The page does not name those textbooks, so their existence and availability are unverified.
\begin{itemize}
\item Omniglot, Yerukula: \url{https://www.omniglot.com/writing/yerukula.htm}
\end{itemize}

\section{Linguistic and Reference Materials}
\begin{itemize}
\item \emph{A Descriptive Grammar of Yerukala, An Endangered South Dravidian Language}, a 2024 PhD thesis at Dravidian University (Bharath Kumar; Mala Sambramalla): \url{https://shodhganga.inflibnet.ac.in/handle/10603/549629}
\item Wikipedia, Yerukala language (reports 27--45\% lexical similarity with Tamil and 53--81\% with Ravula and Irula): \url{https://en.wikipedia.org/wiki/Yerukala_language}
\item Burrow and Emeneau, \emph{A Dravidian Etymological Dictionary}. This is a family-wide comparative reference. The search showed only its front matter, so it is not confirmed to include Yerukula data: \url{https://dsal.uchicago.edu/dictionaries/burrow/new_frontmatter.html}
\item Lexvo entry for \texttt{yeu}: \url{https://finto.fi/lexvo/en/page/yeu}
\end{itemize}

\section{Religious and Community Resources}
Searches found no Yerukula scripture or community organisation. The Joshua Project people-group profile is the only community-related listing found:
\begin{itemize}
\item Joshua Project, Yerukula in India: \url{https://joshuaproject.net/people_groups/18342}
\end{itemize}

\section{Digital and Audio Resources}
No audio, video or corpus resources were found in the search results.

\section{Curriculum Development Context}
The language is not known to be taught in schools. The descriptive grammar is the strongest foundation for curriculum work. Materials would have to be built from scratch, in cooperation with the community and with Dravidian University researchers. The script is also unsettled, since at least two have been proposed or used. Related Dravidian languages could supply comparative material, but they are not substitutes.
\end{document}